\section{Datasets and preprocessing}\label{sec:data}
\subsection{METR-LA and PEMS-BAY}
METR-LA comprises 207 sensors and 34,272 five-minute timestamp rows from 1 March to 27 June 2012. PEMS-BAY comprises 325 sensors and 52,116 original rows from 1 January to 30 June 2017. These statements describe the acquired benchmark files, whose provenance is recorded against DCRNN \citep{li2018}; the graph and speed-file hashes are retained separately. PEMS-BAY has a 65-minute timestamp gap. Reindexing inserts twelve wholly invalid five-minute rows, giving 52,128 prepared rows. METR-LA requires no inserted rows. Table~\ref{tab:data} compares the dimensions, periods, missingness convention and sampling policy. Figure~\ref{fig:data-profile} shows training-partition speed distributions and original validity, making the empirical input populations visible without treating distributional differences as a causal explanation for transfer results.

The benchmark convention treats zero-valued real speeds as missing, following the audited source mask. It should not be interpreted as a physical assertion that vehicles never stop. Originally invalid entries cannot become training labels, calibration residuals or evaluation targets after filling. SUMO instead retains physical zero speeds as valid and masks empty detector intervals. Table~\ref{tab:data-statistics} reports original validity and observed-speed summaries computed directly from prepared arrays. These are descriptive data statistics, not estimates from the artificially censored reporting stream.

\subsection{Chronological separation and causal features}
Real-data boundaries are selected on the original timestamps before reindexing and window construction. The split proportions are 60\% training, 10\% tuning, 10\% initial calibration and 20\% testing. METR-LA boundaries are 11 May 2012 09:35, 23 May 07:10 and 4 June 04:45; PEMS-BAY boundaries are 19 April 2017 14:45, 7 May 17:05 and 25 May 19:20. Table~\ref{tab:partitions} provides these timestamps and the simulation episode allocations. A forecast window belongs to a partition only when all its targets remain inside that partition. Interpolation across a future boundary is prohibited.

Every model receives twelve historical bins, corresponding to 60 minutes, and forecasts twelve future bins. Reported endpoints use horizons three, six and twelve, corresponding to 15, 30 and 60 minutes. Missing input values use the latest available past observation, with the training sensor median as fallback when no prior valid measurement exists. Filling changes the predictor input, not original validity. Seven channels comprise standardised filled speed, validity, capped observation age divided by 288, and sine/cosine representations of time of day and day of week. Scaling, sensor medians, graph processing and context normalisation use training data only. During replay these features are reconstructed from released packets rather than from a fully observed test array.

\subsection{Controlled SUMO datasets}
SUMO/TraCI 1.27.1 is used for a controlled 10-km, two-lane mainline with 20 detector stations \citep{lopez2018}. The simulator steps every second; detector values are aggregated every 300 seconds. Each episode lasts eight hours, with a one-hour warm-up. The original prepared corpus contains 120 episodes and 11,520 timestamp rows. Sixty episodes train the predictor, fifteen tune it, fifteen initialise calibration and thirty form the held-out test set. Causal state resets at episode boundaries, and synthetic inter-episode gaps are not interpreted as observed traffic.

Four registered physical families distinguish baseline demand (P0), increased demand (P1), a temporary speed restriction (P2) and their combination (P3). The demand generator mixes 90\% passenger vehicles and 10\% trucks. For peak families, the sampled demand multiplier is 1.4--1.8. Restriction families temporarily lower the speed on mainline edge m15 to 8.33 $\ms$ for a sampled 1,200--2,400 seconds, then restore 27.78 $\ms$. These are simulated physical events, separate from packet-release faults. They enable repeatable stress without claiming to reproduce every real congestion mechanism.

Further episode sets are generated for independent replication, powered confirmation and final asymmetric validation. The latter contains forty fresh episodes, balanced across P0--P3, with generation seeds 70001--70040. The locked asymmetric C3 candidate is evaluated on this fresh set; C4 transfer uses the thirty original held-out episodes. The separation is recorded explicitly so favourable fresh validation is not confused with an untouched C4 experiment. Prepared simulation arrays, generation settings and source hashes accompany the public release; raw real benchmark packages require separate acquisition because their exact redistribution rights remain unresolved.

\subsection{Observed targets and audit sampling}
The real evaluation is a fault-focused audit: origins are selected hourly and affected sensors are followed through twelve post-event bins. Simulation analyses use full registered episode sets. Both retain originally valid targets even if outcome feedback is delayed or lost. This distinction prevents scoring from favouring a method merely because difficult labels disappeared from its update stream. Table~\ref{tab:data} therefore states the evaluation frequency alongside dataset dimensions. The dataset profile describes measurement populations; it does not establish national, seasonal or road-type representativeness, which remains a limitation of the study.
